% =========================================================
% Diagrama de secuencia: Alertas de reportes cercanos  (6.3.1)
% CU-34 · RF-29, RF-41 · RN-07, RN-18
% ---------------------------------------------------------
% Estilos y macros: setup/settings.tex, sección G3.
% Primera parte: el usuario activa las alertas. Segunda parte: cuando un
% reporte pasa a validado, un webhook de la base de datos dispara una
% Edge Function que notifica por FCM a los usuarios cercanos.
% =========================================================
\begin{tikzpicture}

  % ---------------- Participantes ----------------
  \sdActor{usr}{0}{Peatón\\usuario}
  \sdContenedor[2.1cm]{app}{2.6}{Aplicación móvil}{Flutter}
  \sdExterno[2.1cm]{dev}{5.2}{Servicios del dispositivo}{Android}
  \sdContenedor[2.1cm]{api}{7.8}{API REST}{Supabase Edge Functions}
  \sdContenedor[2.1cm]{bd}{10.4}{Base de datos}{PostGIS}
  \sdExterno[2.1cm]{push}{13.0}{Servicio push}{Firebase Cloud Messaging}
  \sdIniciar

  % ---------------- Activación (RF-41) ----------------
  \sdMensaje[3]{usr}{app}{activa las alertas\\de reportes\\cercanos}
  \sdMensaje[3]{app}{dev}{solicita ubicación\\en segundo plano\\y notificaciones}
  \sdMensaje{dev}{usr}{muestra los diálogos}
  \sdMensaje{usr}{dev}{concede los permisos}
  \sdRespuesta[2]{dev}{app}{permisos\\concedidos}
  \sdMensaje{app}{push}{solicita el token del dispositivo}
  \sdRespuesta{push}{app}{token de FCM}
  \sdMensaje[2]{app}{api}{registra el token\\y la preferencia}
  \sdMensaje{api}{bd}{guarda el token}

  \sdAbrir{fondo}{loop}{[con la aplicación en segundo plano]}{app}{bd}
    \sdMensaje{dev}{app}{posición}
    \sdMensaje[3]{app}{api}{actualiza la\\última ubicación\\conocida}
    \sdMensaje{api}{bd}{guarda la ubicación}
  \sdCerrar{fondo}

  % ---------------- Notificación de un reporte validado ----------------
  \sdMensaje[2]{bd}{api}{webhook: reporte\\validado}
  \sdMensaje[2]{api}{bd}{busca usuarios\\a 1.5 km o menos}
  \sdRespuesta[2]{bd}{api}{tokens de los\\dispositivos}
  \sdAbrir{envio}{loop}{[por cada dispositivo cercano]}{usr}{push}
    \sdMensaje[2]{api}{push}{envía la notificación\\con prioridad alta}
    \sdMensaje{push}{dev}{entrega la notificación}
    \sdMensaje[2]{dev}{usr}{muestra la alerta del\\reporte cercano}
  \sdCerrar{envio}

  \sdTerminar
\end{tikzpicture}
